% ej_coord_R2_03.tex
%
% Copyright (C) 2022--2026 José A. Navarro Ramón <janr.devel@gmail.com>
% Licencia del código GPLv2
% Licencia Creative Commons Recognition Non-Commercial Share-alike.
% (CC-BY-NC-SA)

\section{Ejercicio 3}
Conocido el gradiente de un campo escalar $\phi=\phi(x,y)$, en coordenadas cartesianas ortogonales
\[
  \vvv{\nabla}\phi
  = \dfrac{\partial\phi}{\partial x}\,\uvec{\i}
  + \dfrac{\partial\phi}{\partial y}\,\uvec{\j}
\]
\noindent
a) Calcula sus componentes contravariantes en un sistema de coordenadas oblicuas, cuyos
ejes forman un ángulo $\alpha$ entre sí.\\
b) Calcula sus componentes covariantes en el sistema de coordenadas oblicuas.\\

\noindent
\textbf{Solución}

\noindent
a) Componentes contravariantes en función de $\partial\phi/\partial x$ y
$\partial\phi/\partial y$.

Obtenemos los vectores $\uvec{\i}$ y $\uvec{\j}$ en función de los vectores
unitarios de la base natural $B=(\xhat{e}_1, \xhat{e}_2)$, en coordenadas oblicuas que
forman un ángulo $\alpha$, utilizando la matriz inversa de cambio de base
$\mmm{M}^{-1}$, ver \eqref{eq:R2-oblic_Minv}
\[
  \begin{pmatrix}
    \uvec{\i} & \uvec{\j}
  \end{pmatrix}
  =
  \begin{pmatrix}
    \xhat{e}_1 & \xhat{e}_2
  \end{pmatrix}
  \,
  \begin{pmatrix}
    1 & -\cot\alpha\\
    0 & \phantom{-}\csc\alpha
  \end{pmatrix}
  =
  \begin{pmatrix}
    \xhat{e}_1 & -\cot\alpha\,\xhat{e}_1 + \csc\alpha\,\xhat{e}_2
  \end{pmatrix}
\]

Así obtenemos que
\begin{align*}
  \uvec{\i} &= \xhat{e}_1\\
  \uvec{\j} &= -\cot\alpha\,\xhat{e}_1 + \csc\alpha\,\xhat{e}_2
\end{align*}

Ahora sustituimos en la expresión del gradiente en cartesianas ortogonales
\begin{align*}
  \vvv{\nabla}\phi
  &=
    \dfrac{\partial\phi}{\partial x}\,\uvec{\i}
    + \dfrac{\partial\phi}{\partial y}\,\uvec{\j}
    = \dfrac{\partial\phi}{\partial x}\,\xhat{e}_1
    + \dfrac{\partial\phi}{\partial y}
    \left(-\cot\alpha\,\xhat{e}_1+\csc\alpha\,\xhat{e}_2\right)\\
  &=
    \dfrac{\partial\phi}{\partial x}\,\xhat{e}_1
    - \cot\alpha\dfrac{\partial\phi}{\partial y}\,\xhat{e}_1
    + \csc\alpha\dfrac{\partial\phi}{\partial y}\,\xhat{e}_2\\
  &=
    \left(%
    \dfrac{\partial\phi}{\partial x}-\cot\alpha\dfrac{\partial\phi}{\partial y}
    \right)\,\xhat{e}_1
    + \csc\alpha\,\dfrac{\partial\phi}{\partial y}\,\xhat{e}_2
    = (\nabla\phi)^1\,\xhat{e}_1 + (\nabla\phi)^2\,\xhat{e}_2
\end{align*}
Así
\begin{align*}
  (\nabla\phi)^1
  &=
    \dfrac{\partial\phi}{\partial x^1}
    = \dfrac{\partial\phi}{\partial x}-\cot\alpha\dfrac{\partial\phi}{\partial y}\\
  (\nabla\phi)^2
  &=
    \dfrac{\partial\phi}{\partial x^2}
    = \csc\alpha\,\dfrac{\partial\phi}{\partial y}
\end{align*}
  
\noindent
b) Componentes covariantes en función de $\partial\phi/\partial x$ y
$\partial\phi/\partial y$.

Bajamos índices utilizando el tensor métrico
\[
  (\nabla\phi)_i = g_{ij}\,(\nabla\phi)^j
\]

\begin{align*}
  (\nabla\phi)_1
  &=
    g_{11}(\nabla\phi)^1 + g_{12}(\nabla\phi)^2
    = (\nabla\phi)^1 + \cos\alpha (\nabla\phi)^2\\
  &=
    \dfrac{\partial\phi}{\partial x}-\cot\alpha\dfrac{\partial\phi}{\partial y}
    + \cos\alpha\csc\alpha\dfrac{\partial\phi}{\partial y}
    = \dfrac{\partial\phi}{\partial x}
    - \cancelout{\cot\alpha\dfrac{\partial\phi}{\partial y}}
    + \cancelout{\cot\alpha\dfrac{\partial\phi}{\partial y}}
    =
    \dfrac{\partial\phi}{\partial x}
\end{align*}

\begin{align*}
  (\nabla\phi)_2
  &=
    g_{21}(\nabla\phi)^1 + g_{22}(\nabla\phi)^2
    = \cos\alpha (\nabla\phi)^1 + (\nabla\phi)^2\\
  &=
    \cos\alpha
    \left(%
    \dfrac{\partial\phi}{\partial x}-\cot\alpha\dfrac{\partial\phi}{\partial y}
    \right)
    + \csc\alpha\,\dfrac{\partial\phi}{\partial y}\\
  &= \cos\alpha\,\dfrac{\partial\phi}{\partial x}
    - \dfrac{\cos^2\alpha}{\sin\alpha}\,\dfrac{\partial\phi}{\partial y}
    + \dfrac{1}{\sin\alpha}\,\dfrac{\partial\phi}{\partial y}
    = \cos\alpha\,\dfrac{\partial\phi}{\partial x}
    + \sin\alpha\,\dfrac{\partial\phi}{\partial y}
\end{align*}

Así
\begin{align*}
  (\nabla\phi)_1
  &=
    \dfrac{\partial\phi}{\partial x_1}
    = \dfrac{\partial\phi}{\partial x}\\
  (\nabla\phi)_2
  &=
    \dfrac{\partial\phi}{\partial x_2}
    = \cos\alpha\,\dfrac{\partial\phi}{\partial x}
    + \sin\alpha\,\dfrac{\partial\phi}{\partial y}
\end{align*}


%%% Local Variables:
%%% mode: latex
%%% TeX-engine: luatex
%%% TeX-master: "../retazosmatematicas.tex"
%%% End:
